% !TeX root = ../main.tex
\section{硬件实现与实验评估}
\label{sec:evaluation}

\subsection{实验方法}
\label{sec:eval_setup}

基线及五种变体均由 Synopsys Design Compiler 在 12\,nm、TT、0.7\,V、85$^\circ$C 下综合，采用 1\,GHz 时钟和集成时钟门控，保持八级流水、每拍一笔请求，并满足综合后的时序与电气约束。面积包括组合逻辑、流水寄存器及门控单元。功耗输入取自 Qwen3-8B 第一层 MLP down-projection~\cite{qwen3technicalreport}，分别编码为 BF16、MXFP8 E4M3 和 MXFP4 E2M1。VCS 对综合后网表连续发射 1024 笔独立请求，在中间 512 个周期采集原生 SAIF；Power Compiler 基于完整活动注释估计内部、切换和泄漏功耗~\cite{synopsysActivityPower}。RTL 与门级结果逐笔核对。以下均为综合后、CTS/布线前估计。

\subsection{基线成本与同类工作比较}
\label{sec:reuse_evaluation}
\label{sec:area_breakdown}
\label{sec:related_dp_evaluation}

图~\ref{fig:baseline_breakdown} 显示，基线乘法占面积的 24.15\%，累加与对齐合计占 38.91\%，解包另占 16.02\%，说明专用化需同时裁剪算术和格式处理。H/Q 的对齐功耗占比分别为 20.51\%/20.74\%，E 绕过局部对齐后降至 13.31\%，但解包占比升至 21.31\%。表~\ref{tab:related_dp} 仅列与本文重叠精度的能效和面效，分别按 $2K_pf/P_p$ 和 $2K_pf/A$ 计算；不同工艺、累加精度与活动条件下的数值不作同工况排名。

\begin{figure}[t]
\centering
\subfloat[Area\label{fig:baseline_area}]{\includegraphics[width=0.60\columnwidth]{v19_six/v19_dp_top_area_functions.png}}\par
\subfloat[Power\label{fig:baseline_power}]{\includegraphics[width=\columnwidth]{v19_six/v19_dp_top_power_functions.png}}
\caption{基线按功能归集的面积及总功耗。各类别包含相关寄存；H/Q/E 分别采用 BF16/MXFP8/MXFP4，功耗包含动态与泄漏分量。}
\label{fig:baseline_breakdown}
\end{figure}

% Sources and arithmetic conversions: references/notes/dp_comparison.md.
\begin{table*}[t]
\centering
\caption{Implementation conditions and precision-specific efficiency. Columns retain each work's reported technology and arithmetic boundary.}
\label{tab:related_dp}
\begingroup\footnotesize\setlength{\tabcolsep}{4pt}
\begin{tabular*}{\textwidth}{@{\extracolsep{\fill}}l*{6}{c}@{}}\toprule
Metric & Ten-Four~\cite{rout2026tenfourv2} & MPICC~\cite{huangMPICCMultiplePrecisionInterCombined2025} & MXDOTP~\cite{islamoglu2025mxdotp} & FPnew~\cite{mach2020fpnew} & TransDot~\cite{transdot2026} & This work \\\midrule
Structure & Fused DP & $n$-term MAC & Scaled DP & Vector FMA & FMA / DP & Scaled DP \\
Evaluation & Synthesis & Synthesis & Post-layout & Silicon & Post-layout & Synthesis \\
Technology (nm) & 7 (predictive) & 28 & 12 & 22FDX & 12 & 12 \\
Voltage (V) & 0.7 & 0.9 & 0.8 & 0.8 & 0.8 & 0.7 \\
Frequency (GHz) & 1.5$^a$ & 1.43 & 1.09 & 0.923 & 1.0 & 1.0 \\
Area ($\mu$m$^2$) & 1959.86 & 2295 & 3150 & 49\,000$^b$ & --- & 16\,165.12 \\
Accumulation & FP32/INT32 & FP12/INT8 & FP32 & Mode-specific & FP32$^c$ & FP32/INT32 \\
\midrule
\multicolumn{7}{l}{Energy efficiency (GFLOP/s/W)} \\
\quad FP16 & --- & --- & --- & 497.67 & 555.56 & --- \\
\quad BF16 & --- & --- & --- & 581.96 & --- & 2087.65 \\
\quad FP8 / MXFP8 & --- & 1948 & 2035 & 1244.78 & 1190.48 & 3999.50 \\
\quad FP4 / MXFP4 & --- & 5931 & --- & --- & 2439.02 & 11\,432.11 \\
\quad FP12 (emulated) & --- & 865.1 & --- & --- & --- & --- \\
\midrule
\multicolumn{7}{l}{Area efficiency (GFLOP/s/mm$^2$)} \\
\quad FP16 & --- & --- & --- & 151.43 & --- & --- \\
\quad BF16 & --- & --- & --- & 151.43 & --- & 1979.57 \\
\quad FP8 / MXFP8 & --- & 1246 & 5523.81 & 302.65 & --- & 3959.14 \\
\quad FP4 / MXFP4 & --- & 4985 & --- & --- & --- & 7918.28 \\
\quad FP12 (emulated) & --- & 415.4 & --- & --- & --- & --- \\
\bottomrule\end{tabular*}
\par\smallskip
\parbox{\textwidth}{\raggedright
---: unavailable here or inapplicable; it does not mean zero. One MAC/FMA is counted as two FLOPs; an $n$-term MAC is counted as $2n$ FLOPs.
$^a$Ten-Four targets 1.5 GHz and reports $F_{\max}=1.571$ GHz and total power 6.28 mW; no format-specific efficiency is inferred from that aggregate power.
$^b$FPnew area and nominal SIMD energy include the entire FPU; its FP8 encoding differs from MXFP8.
$^c$TransDot entries use dot products accumulated into FP32; energy efficiency is converted from reported pJ/FLOP.
MXDOTP and this work use MXFP8; this work's FP4 entries use MXFP4. FPnew/MXDOTP area efficiency is calculated from reported throughput and area. Different process nodes, accumulation formats and activity conditions preclude a direct iso-technology ranking.}
\endgroup
\end{table*}


\subsection{专用化与 K/N 组合成本}
\label{sec:lowp_costs}
\label{sec:kn_evaluation}

图~\ref{fig:six_dp_costs} 比较六种 DP。E64/Q32E64 面积分别为基线的 34.92\%/72.35\%。E96/E128 以 16.69\%/33.60\% 的面积增长换取 50\%/100\% 的 E 算力增长；Q64E128 的 Q/E 算力翻倍，但 Q 功耗达到 28.789\,mW。增加低精度算力并不意味着总功耗按面积同比增长。

\begin{figure}[t]
\centering
\subfloat[Area\label{fig:six_dp_area}]{\includegraphics[width=0.94\columnwidth]{latest6_area.png}}\par
\subfloat[Power\label{fig:six_dp_power}]{\includegraphics[width=0.94\columnwidth]{latest6_power.png}}
\caption{基线及五种变体的面积和总功耗。图中面积倍率以基线面积归一化，功耗倍率均以基线 BF16 功耗归一化；缺失柱表示不支持该精度。}
\label{fig:six_dp_costs}
\end{figure}

K/N 比较采用 DP 成本叠加模型：阵列面积仅计完整 DP 面积之和，功耗仅计活动 DP 的模式功耗，忽略存储、互连及其他逻辑。N 方案中不支持当前精度的附加 DP 假设理想 power gating，动态与泄漏功耗均为零，不计关断与唤醒开销；其物理面积仍保留。表~\ref{tab:kn_costs} 按基线 DP 数归一化，使两方案各精度的峰值一致。新增单元参与时，K 的面积低 0.65\%--2.40\%，功耗低 2.02\%--3.97\%；H 模式下，理想关断的 N 方案功耗更低。独立 E64 与内部 E 构件的局部流水及规约实现不同，因此差额包含实现差异，并非仅改变组合位置的消融结果。

% Generated by paper/data/kn_comparison/calculate.py from retained DP measurements.
\begin{table}[t]
\centering
\caption{Modeled costs per baseline DP (K / N). Area: $10^3\,\mu$m$^2$; power: mW}
\label{tab:kn_costs}
\begingroup\footnotesize\setlength{\tabcolsep}{2pt}
\begin{tabular*}{\columnwidth}{@{\extracolsep{\fill}}lcccc@{}}\toprule
Ratio & Area (K / N) & $P_H$ (K / N) & $P_Q$ (K / N) & $P_E$ (K / N) \\\midrule
$1:2:6$ & 18.863/18.987 & 15.637/15.328 & 16.354/16.002 & 14.417/14.715 \\
$1:2:8$ & 21.596/21.809 & 15.893/15.328 & 16.620/16.002 & 17.595/18.233 \\
$1:4:8$ & 27.191/27.860 & 15.744/15.328 & 28.789/29.582 & 20.563/21.414 \\
\bottomrule
\end{tabular*}\endgroup
\end{table}


\subsection{阵列复用与累加带宽}

采用稳态行列复用模型：A 沿行共享、B 沿列共享，广播或脉动供数均按理想复用计，忽略填充与排空。对某精度扩展 $r$ 倍算力时，K/N 方案分别为 $m\times n\times rk$ 和 $m\times rn\times k$。若原生输入向量（含尺度）占 $V$ bit，完整缩放块扩展后占 $rV$ bit，累加位宽为 $c$ bit（FP32 取 32），所分析存储接口每 $\tau$ 次更新交换一次 C/D，则
\begin{equation}
\begin{aligned}
B_K &= f[(m+n)rV+2mnc/\tau],\\
B_N &= f[(m+rn)V+2rmnc/\tau],\\
B_N-B_K &= f(r-1)m(2nc/\tau-V).
\end{aligned}
\label{eq:kn_bandwidth_difference}
\end{equation}
其中带宽单位为 bit/s，两方案取相同 $\tau$。N 的每个累加上下文由 $mn$ 个结果增至 $rmn$ 个，存储容量与 C/D 端口需求均增至 $r$ 倍，同时通过行内复用减少 A 供数。当 $2nc/\tau>V$ 时，K 的总接口带宽更小；长期驻留部分和或输入向量较宽时，N 的输入复用优势更明显。该式比较稳态接口压力，并不表示完整 GEMM 的最终输出量增加。基于 DP 成本与累加状态开销，后续负载分析采用 K 扩展。
